\section{Introduction}

State-space models are used when a latent Markov process is observed with noise and a static parameter governs the transition and the observation. The observed-data likelihood is an integral over the latent path, and that integral is rarely available in closed form. Sequential Monte Carlo squared (SMC2) propagates a cloud of parameter particles, reweights them by a particle-filter estimate of each likelihood increment, and moves them with a particle Metropolis--Hastings kernel when the effective sample size falls \citep{chopin2012}. The variance of that estimate grows with the length of the series and, more sharply, with the observation dimension.

The ensemble Kalman filter replaces the forecast distribution by a Gaussian distribution whose mean and covariance are the sample moments of an ensemble, and it updates the ensemble with a Kalman shift. The likelihood contribution is then the Gaussian marginal of that forecast. It is not unbiased for the true likelihood. \citet{golightly2026} place this contribution inside the SMC2 recursion. The resulting nested ensemble Kalman filter targets the parameter posterior induced by the ensemble Kalman likelihood, which we denote \(\hat{\pi}_{\mathrm{enkf}}\), rather than the exact posterior \(\pi_t\).

For a nonlinear or non-Gaussian observation the same construction replaces the observation density by a linear Gaussian approximation before the ensemble update. The plain filter still targets \(\hat{\pi}_{\mathrm{enkf}}\), now under that approximation. A Rao-Blackwellised weight puts the true observation density back into the increment, while the forecast remains Gaussian, and targets a different posterior, which the source paper denotes \(\tilde{\pi}_{\mathrm{enkf}}\). Vanilla SMC2, with a bootstrap proposal, targets \(\pi_t\). Separately, the nearest-neighbour surrogate that \citet{golightly2026} use to screen a Metropolis--Hastings proposal can be used as the incremental weight itself. That choice defines a further posterior, which we denote \(\pi_{\mathrm{knn}}\). A screen that is corrected by the ensemble Kalman likelihood leaves \(\hat{\pi}_{\mathrm{enkf}}\) invariant. A weight that is the surrogate does not.

These are two errors, and a comparison that jumps from \(\pi_{\mathrm{knn}}\) to \(\pi_t\) mixes them. The first error is the ensemble Kalman likelihood against the true likelihood, hence \(\hat{\pi}_{\mathrm{enkf}}\) against \(\pi_t\). The second is the nearest-neighbour prediction against the ensemble Kalman likelihood, hence \(\pi_{\mathrm{knn}}\) against \(\hat{\pi}_{\mathrm{enkf}}\). On the Lotka--Volterra diffusion in Section~5.2 of \citet{golightly2026}, the observation is Gaussian with variance equal to the latent prey state. That paper reports that the plain filter shifts the posterior means of \(\theta_2\) and \(\theta_3\) because the observation variance is replaced by a function of the forecast mean, and that the Rao-Blackwellised filter reduces the discrepancy at greater cost. The distance between \(\pi_{\mathrm{knn}}\) and \(\hat{\pi}_{\mathrm{enkf}}\) is not bounded there.

Our contribution is to separate the two errors and to compare the three targets on that nonlinear observation model. (i)~The plain nested filter targets \(\hat{\pi}_{\mathrm{enkf}}\). The nearest-neighbour incremental weight targets \(\pi_{\mathrm{knn}}\), and we do not treat those posteriors as the same object. (ii)~When the true observation density equals the linear Gaussian density inside the filter, the Rao-Blackwellised weight equals the ensemble Kalman weight. The identity is specific to that case. (iii)~On the Lotka--Volterra observation we compare the plain filter, the Rao-Blackwellised filter and bootstrap SMC2, under a reduced design that is stated as such. A bound on the distance between \(\pi_{\mathrm{knn}}\) and \(\hat{\pi}_{\mathrm{enkf}}\) was not obtained. The paper reports the diagnostic that was fixed before any output was seen. A faster neighbour search, a kernel chosen by time, and a finite-difference proposal were examined in an earlier pilot and are not contributions here.

The identity is checked on a scalar Ornstein--Uhlenbeck diffusion with a known Kalman likelihood, at a smaller particle count and a shorter series than Section~5.1 of \citet{golightly2026}. The nonlinear comparison uses one prey series from the Lotka--Volterra diffusion, with fewer parameter particles and fewer algorithm replicates than Section~5.2, and without the particle Markov chain Monte Carlo gold standard used there. No epidemic series is analysed. \citet{temfack2026} already study a nested ensemble scheme for epidemic incidence.

The remainder of this paper is organised as follows. Section~\ref{sec:related} describes the neighbouring algorithms and which posterior each one targets. Section~\ref{sec:methods} separates the two errors and records the identity. Section~\ref{sec:sim} reports the identity check and the Lotka--Volterra comparison, and Section~\ref{sec:disc} discusses what those runs support and what remains open.
